\chapter{LCS con un numero pari di simboli rossi}

\section{Definizione del problema}

\subsection{Istanza}
Siano date due sequenze di elementi sull'alfabeto $\Sigma$:
\begin{align*}
    X &= \langle x_1, x_2, \dots, x_m \rangle \\
    Y &= \langle y_1, y_2, \dots, y_n \rangle
\end{align*}
e una funzione di colorazione $\text{col} : \Sigma \to \{\text{rosso}, \text{blu}, \text{nero}\}$.

\subsection{Soluzione}
Una più lunga sottosequenza comune tra $X$ e $Y$ la quale contiene un numero pari di simboli di colore rosso.

\section{Definizione del sottoproblema e del coefficiente}
Per ogni $i \in \{0, \dots, m\}$, per ogni $j \in \{0, \dots, n\}$ e per ogni $p \in \{0, 1\}$, siano $X_i$ e $Y_j$ i prefissi di lunghezza $i$ e $j$ di $X$ e $Y$, e sia $p$ un indice booleano di parità ($k \bmod 2 = p$, dove $k$ è il numero di elementi rossi nella sottosequenza).

\medskip
\noindent
\textbf{Definizione del coefficiente $c_{i,j}^p$:} Per ogni $i \in \{0, \dots, m\}$, $j \in \{0, \dots, n\}$ e $p \in \{0, 1\}$, $c_{i,j}^p$ è la lunghezza della più lunga sottosequenza comune tra $X_i$ e $Y_j$ la quale contiene:
\begin{itemize}
    \item un numero \textbf{pari} di elementi rossi, se $p = 0$;
    \item un numero \textbf{dispari} di elementi rossi, se $p = 1$.
\end{itemize}
Qualora non esista alcuna sottosequenza comune con la parità richiesta, si pone $c_{i,j}^p = -\infty$.

\section{Equazioni di ricorrenza}

\subsection{Caso base}
\begin{align*}
    (i = 0 \lor j = 0) \land p = 0 &\implies c_{i,j}^0 = 0 \\
    (i = 0 \lor j = 0) \land p = 1 &\implies c_{i,j}^1 = -\infty
\end{align*}
\textbf{Giustificazione:} Se almeno uno dei due prefissi è vuoto, l'unica sottosequenza comune è $\langle \rangle$, che contiene $0$ elementi rossi. Poiché $0$ è un numero pari, per $p = 0$ la lunghezza ottima è $0$; non potendo invece avere un numero dispari di elementi rossi nel prefisso vuoto, per $p = 1$ lo stato è inammissibile ($-\infty$).

\subsection{Passo ricorsivo}
Per $i > 0 \land j > 0$ e per ogni $p \in \{0, 1\}$:
\begin{itemize}
    \item \textbf{Caso $x_i \neq y_j$:}
    \[
        c_{i,j}^p = \max \{ c_{i-1, j}^p, \, c_{i, j-1}^p \}
    \]
    \textbf{Giustificazione:} Non essendovi match tra $x_i$ e $y_j$, la parità richiesta $p$ rimane invariata e si prende il massimo tra i sottoproblemi $(i-1, j, p)$ e $(i, j-1, p)$.

    \item \textbf{Caso $x_i = y_j$ con $\text{col}(x_i) \neq \text{rosso}$:}
    \[
        c_{i,j}^p = c_{i-1, j-1}^p + 1
    \]
    \textbf{Giustificazione:} Un elemento non rosso lascia inalterata la parità del numero di simboli rossi e incrementa di $1$ la lunghezza della sottosequenza ottima proveniente da $(i-1, j-1, p)$.

    \item \textbf{Caso $x_i = y_j$ con $\text{col}(x_i) = \text{rosso}$:}
    \begin{itemize}
        \item Se $p = 0$:
        \[
            c_{i,j}^0 = \max \{ c_{i-1, j-1}^1 + 1, \, c_{i-1, j-1}^0 \}
        \]
        \textbf{Giustificazione:} Per ottenere un numero pari di rossi ($p=0$), possiamo includere il rosso corrente $x_i = y_j$ agganciandolo a una sottosequenza con numero dispari di rossi ($c_{i-1, j-1}^1 + 1$), oppure scartarlo mantenendo la migliore sottosequenza già pari ($c_{i-1, j-1}^0$).
        \item Se $p = 1$:
        \[
            c_{i,j}^1 = \max \{ c_{i-1, j-1}^0 + 1, \, c_{i-1, j-1}^1 \} = c_{i-1, j-1}^0 + 1
        \]
        \textbf{Giustificazione:} Per ottenere un numero dispari di rossi ($p=1$), includendo il rosso corrente si passa da una sottosequenza con numero pari di rossi ($c_{i-1, j-1}^0 + 1$). Si noti che vale sempre $c_{i-1, j-1}^0 + 1 \ge c_{i-1, j-1}^1$, poiché rimuovendo un elemento rosso da una qualsiasi sottosequenza con numero dispari di rossi si ottiene una sottosequenza valida con numero pari di rossi di lunghezza $c_{i-1, j-1}^1 - 1 \le c_{i-1, j-1}^0$.
    \end{itemize}
\end{itemize}

\subsection{Soluzione del problema principale}
La lunghezza della più lunga sottosequenza comune tra $X$ e $Y$ con un numero pari di simboli rossi è:
\[
    \text{Soluzione} = c_{m, n}^0
\]

\section{Algoritmo Bottom-Up}
\begin{algorithm}[htbp]
\caption{Calcolo Bottom-Up per LCS con un numero pari di simboli rossi}
\begin{algorithmic}[1]
\Procedure{LCSRP}{$X, Y, \text{col}$}
    \State $m \gets |X|$
    \State $n \gets |Y|$
    \State $c \gets \text{nuova matrice } c[0 \dots m, \, 0 \dots n, \, 0 \dots 1]$
    \Comment{Caso base}
    \For{$i = 0 \textbf{ to } m$}
        \State $c[i, 0, 0] \gets 0$
        \State $c[i, 0, 1] \gets -\infty$
    \EndFor
    \For{$j = 0 \textbf{ to } n$}
        \State $c[0, j, 0] \gets 0$
        \State $c[0, j, 1] \gets -\infty$
    \EndFor
    \Comment{Passo ricorsivo}
    \For{$i = 1 \textbf{ to } m$}
        \For{$j = 1 \textbf{ to } n$}
            \For{$p = 0 \textbf{ to } 1$}
                \If{$x_i \neq y_j$}
                    \State $c[i, j, p] \gets \max \{ c[i-1, j, p], \, c[i, j-1, p] \}$
                \Else
                    \If{$\text{col}(x_i) \neq \text{rosso}$}
                        \State $c[i, j, p] \gets c[i-1, j-1, p] + 1$
                    \Else
                        \If{$p == 0$}
                            \State $c[i, j, 0] \gets \max \{ c[i-1, j-1, 1] + 1, \, c[i-1, j-1, 0] \}$
                        \Else
                            \State $c[i, j, 1] \gets c[i-1, j-1, 0] + 1$
                        \EndIf
                    \EndIf
                \EndIf
            \EndFor
        \EndFor
    \EndFor
    \State \Return $c$
\EndProcedure
\end{algorithmic}
\end{algorithm}

\section{Algoritmo di ricostruzione}
\begin{algorithm}[htbp]
\caption{Ricostruzione della sequenza LCS con numero pari di rossi}
\begin{algorithmic}[1]
\Procedure{LCSRPR}{$c, X, Y, i, j, p$}
    \If{$i == 0 \lor j == 0$}
        \State \Return
    \Else
        \If{$x_i \neq y_j$}
            \If{$c[i-1, j, p] \ge c[i, j-1, p]$}
                \State \Call{LCSRPR}{$c, X, Y, i-1, j, p$}
            \Else
                \State \Call{LCSRPR}{$c, X, Y, i, j-1, p$}
            \EndIf
        \Else
            \If{$\text{col}(x_i) \neq \text{rosso}$}
                \State \Call{LCSRPR}{$c, X, Y, i-1, j-1, p$}
                \State \textbf{print} $x_i$
            \Else
                \If{$p == 0$}
                    \If{$c[i-1, j-1, 1] + 1 \ge c[i-1, j-1, 0]$}
                        \State \Call{LCSRPR}{$c, X, Y, i-1, j-1, 1$}
                        \State \textbf{print} $x_i$
                    \Else
                        \State \Call{LCSRPR}{$c, X, Y, i-1, j-1, 0$}
                    \EndIf
                \Else
                    \State \Call{LCSRPR}{$c, X, Y, i-1, j-1, 0$}
                    \State \textbf{print} $x_i$
                \EndIf
            \EndIf
        \EndIf
    \EndIf
\EndProcedure
\end{algorithmic}
\end{algorithm}

\noindent
\textbf{Chiamata principale (Driver):} Si invoca $\Call{LCSRPR}{c, X, Y, m, n, 0}$ (notando che una soluzione esiste sempre poiché la sequenza vuota ha $0$ rossi e garantisce $c[m, n, 0] \ge 0$).

\section{Analisi della complessità}
\begin{itemize}
    \item \textbf{Complessità spaziale:} $\Theta(2 \cdot m \cdot n) = \Theta(m \cdot n)$.
    \item \textbf{Complessità temporale (Bottom-Up):} $\Theta(2 \cdot m \cdot n) = \Theta(m \cdot n)$.
    \item \textbf{Complessità temporale (Ricostruzione):} $O(m + n)$.
\end{itemize}